\section*{Problem 2 (10 points)}

Throughout, the assignment's binary unit convention applies, so
$64\text{ KB} = 2^{16}$ Bytes and $2\text{ GB} = 2^{31}$ Bytes.

\subsection*{Question A (2 points)}

The page offset must address every byte within a page. A page holds
$64\text{ KB} = 2^{16}$ Bytes, so the offset is $16$ bits, and the remaining
$30 - 16 = 14$ bits of the $30$-bit virtual address form the virtual page number.

\begin{table}[H]
  \centering
  \caption{virtual address components}
  \label{p2:tab:vaddr}
  \begin{tabular}{ccc}
    \toprule
    Page offset / bits & Virtual page memory (VPN) / bits & Total / bits \\
    \midrule
    16 & 14 & 30 \\
    \bottomrule
  \end{tabular}
\end{table}

\subsection*{Question B (4 points)}

A single-level page table holds one entry per virtual page, so it is indexed by the
$14$-bit VPN from Question~A; the physical memory size does not determine the number
of entries.

\textbf{Number of PTEs:} $2^{14} = 16{,}384$ entries.

\textbf{Physical memory for the page table:}
$2^{14} \times 4\text{ Bytes} = 2^{16}\text{ Bytes} = 64\text{ KB}$.

\subsection*{Question C (4 points)}

\textbf{Pros.}
\begin{enumerate}
  \item \emph{Smaller page table.} A larger page size widens the offset field and
    narrows the VPN by the same number of bits, so the entry count falls
    geometrically. Doubling the page size in Question~B to $128$ KB would leave a
    $13$-bit VPN, halving the table from $64$ KB to $32$ KB.
  \item \emph{Greater TLB reach.} Each TLB entry maps a whole page, so with larger
    pages a fixed number of TLB entries covers proportionally more of the address
    space. That lowers the TLB miss rate for programs with large working sets, and
    each avoided miss avoids a page-table walk.
\end{enumerate}

\textbf{Cons.}
\begin{enumerate}
  \item \emph{More internal fragmentation.} Memory is allocated in whole pages, so
    the unused remainder of a partially filled page is wasted. The waste per
    incompletely filled region grows with the page size, which hurts most for many
    small objects or sparse address spaces.
  \item \emph{Costlier page faults.} A fault transfers an entire page from disk, so
    a larger page means a longer transfer and more memory traffic. If the program
    only touches a small part of that page, the extra bytes fetched are wasted I/O
    and the fault latency rises.
\end{enumerate}
